\paragraph{First-order equivariant class on left-invariant backgrounds with curvature and torsion: exact solutions and the linearised equations (conditional; not closure).}
Keep the class of the preceding paragraph on a four-dimensional Lie group, with left-invariant frame \(E_a\), \([E_a,E_b]=\mathsf c_{ab}{}^dE_d\) and \(\mathbb D_a=u_a{}^bE_b\) plus the connection term of the class, but give the connection constant components \(w_a{}^{cd}\) in this frame, normalised so that the first-order part of \(\mathbb D_a\) at \(u=1\) is \(E_a+\tfrac12w_a{}^{cd}\mathcal S_{cd}\), the sign of the fluctuation \(\tfrac12\omega_a{}^{cd}\mathcal S_{cd}\) below; in the coordinate form of the paragraph on local \(\mathrm{SO}(3,1)\) inside the unitary group this is \(e_a{}^\mu\omega_\mu{}^{cd}=-w_a{}^{cd}\). The section values \(\mathsf T\) and \(\mathsf R\) are then constants, polynomial of degree at most two in \(\mathsf c\) and \(w\); they agree with the relations computed for the vector fields on the bundle at random Lorentz \(u\) (numerical, residual below \(10^{-13}\)), and the first two brackets of the matrix equation become polynomial equations in \(\mathsf c\) and \(w\). With \(\eta=\mathrm{diag}(-1,1,1,1)\), let \(B^{cd}\) be the boost bivector with \(B^{02}=1\), let \(N^{cd}\) and \(N'^{cd}\) be the null rotations with \(N^{01}=N^{12}=1\) and \(N'^{03}=-N'^{23}=1\), which fix the null direction \(n^a=(1,0,-1,0)\), and let \(k_a=(1,0,1,0)\), so that \(k^a=-n^a\); components not listed vanish apart from antisymmetry. Two families solve both brackets exactly, for symbolic parameters, with \(\mathsf R\neq0\) and \(\mathsf T\neq0\):
\begin{itemize}
\item[(i)] the flat frame, \(\mathsf c=0\), with \(w_a{}^{cd}=s\,k_aB^{cd}+m_aN^{cd}\), \(m_a=(0,m_1,0,m_3)\) and \((m_1-s)^2-m_3^2=s^2\). In this ansatz the second bracket vanishes identically and the first is proportional to \(m_1^2-2m_1s-m_3^2\). Here \(\mathsf R_{ab}{}^{cd}\) has rank one with values on the line of \(N\), \(\mathsf R\neq0\) for \(s\neq0\) unless \(m=0\), \(\mathsf R_{abcd}\mathsf R^{abcd}=\mathsf R_{ab}{}^{ab}=0\), \(\mathsf T_{abc}\mathsf T^{abc}=0\), and \(\mathsf T_{ac}{}^c=(s+m_1)k_a\), which vanishes at \(m=(0,-s,0,\pm\sqrt3\,s)\). The frame metric is that of Minkowski space. For constant section values the order-zero field \(\mathsf t_a=\tfrac12\mathsf T_{ac}{}^c\) of the class gives \(\mathsf F_{ab}=-\mathsf c_{ab}{}^c\mathsf t_c\) (by hand), which is zero here, so the third bracket also holds;
\item[(ii)] time times the Heisenberg group of the preceding paragraph, \([E_3,E_1]=\kappa E_2\), with \(w_a{}^{cd}=k_a(\zeta N^{cd}+xN'^{cd})\) and \((\zeta,x)\neq(0,0)\). The torsion is trace-free, so \(\mathsf t=0\) and the third bracket holds, \(\mathsf R_{13}{}^{cd}=\kappa(\zeta N^{cd}+xN'^{cd})\) is the only nonzero pair, \(\mathsf R_{abcd}\mathsf R^{abcd}=\mathsf R_{ab}{}^{ab}=0\) and \(\mathsf T_{abc}\mathsf T^{abc}=2\kappa^2\).
\end{itemize}
In both families the Ricci tensor \(\mathsf R_{ab}{}^{ae}\) of the connection is nonzero and is cancelled in the first bracket by the torsion terms. On the algebra of (ii) the wider ansatz \(w_a{}^{cd}=k_a(aB^{cd}+\zeta N^{cd}+xN'^{cd})\) solves all three brackets for all real \(\kappa,a,\zeta,x\): \(\mathsf T_{ac}{}^c=a\,k_a\), the only nonzero pair of \(\mathsf F\) is \(\mathsf F_{13}=\tfrac12a\kappa\), the only nonzero pair of \(\mathsf R\) is \(\mathsf R_{13}{}^{cd}=\kappa(aB^{cd}+\zeta N^{cd}+xN'^{cd})\), and \(\mathsf R_{abcd}\mathsf R^{abcd}=-4a^2\kappa^2\), \(\mathsf R_{ab}{}^{ab}=0\), \(\mathsf T_{abc}\mathsf T^{abc}=2\kappa^2\) (exact for the section values; for the operators on the bundle at random Lorentz \(u\), including the order-zero part, residual below \(10^{-15}\)). For \(a\neq0\) it is an exact solution with \(\mathsf R_{abcd}\mathsf R^{abcd}\neq0\) on a unimodular algebra. In a numerical survey with \(\mathsf c=0\), 56 of 60 least-squares starts converged to solutions with \(\mathsf R\neq0\) for \(\eta\), all with \(\mathsf R_{abcd}\mathsf R^{abcd}=0\) to \(10^{-14}\), and none converged for the Euclidean contraction; with \(\mathsf c\) free, solutions of the first two brackets with \(\mathsf R_{abcd}\mathsf R^{abcd}\neq0\) were found on algebras with a two-dimensional abelian derived algebra, not all unimodular, for which the order-zero field and the third bracket are not examined (numerical, not a classification).

On family (i) the coefficients are constant, so the translations \(\partial_I\) commute with all \(\mathbb D_a\) and \(\mathcal S_{cd}\), and at \(u=1\) one has \(\partial_I=\mathbb D_I-\tfrac12w_I{}^{cd}\mathcal S_{cd}\). With \(\mathbb D_a\to\mathbb D_a-h_a{}^c\mathbb D_c+\tfrac12\omega_a{}^{cd}\mathcal S_{cd}\) as before, the adjoint compatibility of the paragraph on the pairing of the action then gives, for each \(I\) and arbitrary functions \(f_a\), that the fluctuation \(h_a{}^c=-f_a\delta^c{}_I\), \(\omega_a{}^{cd}=-f_aw_I{}^{cd}\) has linearised brackets
\[
\mathcal E^{(1)}{}_b{}^e=\mathcal Q_b[f]\,\delta^e{}_I,\qquad
\mathcal E^{(2)}{}_b{}^{de}=-\mathcal Q_b[f]\,w_I{}^{de},
\]
with \(\mathcal Q\) the linearised third bracket on this background, \(\mathsf F_{ab}[A]=\nabla_aA_b-\nabla_bA_a+\mathsf T_{ab}{}^cA_c\) and \(\mathcal Q_b[A]=\nabla^a\mathsf F_{ab}[A]-\mathsf T_{ab}{}^c\mathsf F^a{}_c[A]\), where \(\nabla\) now contains \(w\) (checked at sampled \(q\), residual \(10^{-15}\)). This is the frame reduction of the preceding paragraph at \(\mathsf R\neq0\), with labels that now carry a connection part. The reduction of the connection equations to \(\mathcal Q\) fails, since the linearised curvature contains \(h\): the connection equations depend on pure frame fluctuations (numerical). In the variables \(f_b{}^I=-h_b{}^I\) and \(\omega'_b{}^{cd}=\omega_b{}^{cd}-h_b{}^Iw_I{}^{cd}\), the twenty-four combinations \(\mathcal E^{(2)}{}_b{}^{de}+\mathcal E^{(1)}{}_b{}^Iw_I{}^{de}\) contain \(\omega'\) alone, so the system is again block-triangular, the connection sourcing the frame labels, with a connection block of generic rank 18 that is not six copies of \(\mathcal Q\). The kernel at generic \(q\) is the ten gauge directions, whose diffeomorphism part now has \(\omega_a{}^{de}=\xi^c\mathsf R_{ca}{}^{de}\) (numerical). Along lines \(q=(t,\vec q)\) every characteristic root satisfies \(q^2=2\alpha\,k^aq_a\), that is \((q-\alpha k)^2=0\), for \(\alpha\) in a finite set: at \(s=1\), \(\alpha\in\{\pm\sqrt3,\pm\sqrt3\pm\mathrm i,\pm\sqrt3+2\mathrm i\}\) for \(m=(0,-1,0,\sqrt3)\), with \(\pm\sqrt3\) for \(\mathcal Q\), and \(\alpha\in\{-\tfrac{\mathrm i}2,-\tfrac{3\mathrm i}2,-\tfrac{5\mathrm i}2\}\) for \(m=(0,2,0,0)\), with \(-\tfrac{3\mathrm i}2\) for \(\mathcal Q\). Conversely the forty equations drop rank at random complex points of each of these cones (numerical; the relative smallest singular value is about \(10^{-16}\) at both points, against \(3\cdot10^{-5}\) to \(3\cdot10^{-2}\) at random points off the cones). The background has an exact boost-dilation symmetry: \(\exp(\theta B)\) multiplies \(k_a\) and \(N\) by \(\mathrm e^{-\theta}\) and fixes \(B\) and \(m\), so it maps \(w\) to \(\mathrm e^{-\theta}w\), while \(\mathbb D_a\to\lambda\mathbb D_a\) maps solutions of the matrix equation to solutions. The set of characteristic \(q\) is therefore invariant under \(q\mapsto\mathrm e^{\pm\theta}\Lambda_\theta q\), with \(\Lambda_\theta\) the boost and the sign fixed by the action on covectors, so the set of values of \(q^2\) on it is invariant under multiplication by every positive number, and no value \(q^2=-m^2\) with \(m\neq0\) is isolated. So on family (i) the non-derivative terms, which are nonzero, shift the cones along \(k\) and do not produce a mass: no gauge-invariant gap appears for the Lorentz or torsion modes at this order, and the background supplies no invariant scale for one. The imaginary parts of \(\alpha\) give solutions whose section values contain \(\mathrm e^{\beta k_ay^a}\) with \(\beta\) real, which grow along the null coordinate \(k_ay^a\); since section values are frame components and no norm is used, no statement on stability is made.

These statements are for Equation~\eqref{eq:Smatrix} with \(C_{\mu\nu\lambda}=0\) and no matter term, tree level, the Lorentzian contraction, left-invariant backgrounds with constant connection, the families named, and, for the linearised equations, linear order on family (i) in its character sectors, with the roots numerical. Family (ii), its extension with \(a\neq0\) and the survey solutions with \(\mathsf R_{abcd}\mathsf R^{abcd}\neq0\), for which no scaling symmetry is shown, are not linearised here, and the nonlinear evolution, the regulated trace, the domains of the unbounded operators, loops and a norm are not treated. Nothing here shows that the matrix integral selects this class, these backgrounds or the algebra \(\mathcal A\), or yields the frame and its non-degeneracy, so Assumptions~\ref{ass:frame-lock}, \ref{ass:nondeg}, \ref{ass:lorentzian}, \ref{ass:local-lorentz} and~\ref{ass:EFT-a2} are not discharged, and the frame of these backgrounds is not identified with the emergent metric of Chapter~\ref{ch:vierbein-emergence}; Problems~D, F and~N remain open and PROGRAMME-CONDITIONAL.
